\documentclass[a4paper,11pt,dvipsnames]{article}
\usepackage{amsmath,amsthm,amssymb,amsfonts}
\usepackage{array}
\usepackage{tabularx}
\usepackage[table]{xcolor}
\usepackage[most]{tcolorbox}
\usepackage{graphicx}
\graphicspath{{./img/}}
\usepackage[a4paper,left=1.5cm,right=1.5cm,top=1cm,bottom=1.5cm]{geometry}
% \usepackage[hmargin=1cm,vmargin=1cm]{geometry}
\usepackage{array}
\usepackage{setspace}
\newcommand*\diff{\mathop{}\!\mathrm{d}}
\newcommand*\Diff[1]{\mathop{}\!\mathrm{d^#1}}

\newcommand{\N}{\mathbb{N}}
\newcommand{\Z}{\mathbb{Z}}

% A style for tcolorbox without frame etc. "notcb"
\tcbset{notcb/.style={enhanced jigsaw, colback=white, coltitle=black,sharp corners, boxrule=0pt}}

% Style for the number tables
\tcbset{tabularheader/.style={arc=2pt,notcb,boxsep=0pt,left=0pt,colbacktitle=white,rounded corners=all,enhanced jigsaw,coltitle=black,tabularx={*{6}{Y|}Y},boxrule=1pt,attach boxed title to top left, boxed title style={enhanced jigsaw,sharp corners,left=0pt,boxrule=0pt}}}

\newcolumntype{Y}{>{\centering\arraybackslash}X}

\usepackage[locale=FR, per-mode=fraction, separate-uncertainty=true]{siunitx}
\usepackage{physics}
\usepackage[utf8]{inputenc}
\usepackage[T1]{fontenc}
\usepackage[spanish]{babel}

\usepackage{enumitem}

\usepackage{pgfplots}
% \usepackage{pdfpages}
\usepackage{adjustbox, tikz}
\usepackage{circuitikz}
% \usepackage{gnuplottex}
\usepackage{xcolor}
\usepackage[locale=FR, per-mode=fraction, separate-uncertainty=true]{siunitx}
\usepackage{physics}

\usepackage{ifpdf}
\ifpdf%
  \usepackage{pdftricks}
  \begin{psinputs}
    \usepackage{pst-optic}
%\usepackage{pstricks-add} 	
  \end{psinputs}
\else
  \usepackage{pst-optic}
  \newenvironment{pdfpic}{}{}
\fi
 

\usetikzlibrary{decorations}
\usetikzlibrary{decorations.pathmorphing}
\usetikzlibrary{shapes.geometric}
\usetikzlibrary{arrows}
\usetikzlibrary{arrows.meta}   %edg
\usetikzlibrary{calc}
\pgfplotsset{compat=newest}


\begin{document}

{\Large Resoluciones 2do parcial - 2022}

\begin{enumerate}[leftmargin=*]
  \item Dos cargas puntuales, $q_1 = Q$ y $q_2 = -2Q$, siendo $Q>0$, están ubicadas sobre el eje $x$ como se muestra en la figura. Marcar todas las opciones sugeridas donde la fuerza resultante sobre una carga positiva ubicada sobre el eje $x$, esté dirigida hacia la derecha.
  \begin{center}
    \begin{tikzpicture}[scale=0.8]
      \draw [{latex}-{latex}] (-5,0)--(8,0) node[below] {$x$~(m)};
      \fill [black](0,0) circle(4pt) node[above] {$q_1$};
      \fill [black](3,0) circle(4pt) node[above] {$q_2$};
      % \draw [] (-4.5,0.2) -- (-4.5,-0.2);
      \draw [] (-3,0.2) -- (-3,-0.2) node[below] {$-2$};
      \draw [] (-1.5,0.2) -- (-1.5,-0.2) node[below] {};
      \draw [] (0,0) -- (0,-0.2) node[below] {$0$};
      \draw [] (1.5,0.2) -- (1.5,-0.2) node[below] {};
      \draw [] (3,0) -- (3,-0.2) node[below] {$2$};
      \draw [] (4.5,0.2) -- (4.5,-0.2) node[below] {};
      \draw [] (6,0.2) -- (6,-0.2) node[below] {$4$};
      % \draw [] (-1,0.2) -- (-1,-0.2);
      % \fill [black](5,0) circle(4pt) node[above] {$q_3$};
    \end{tikzpicture}
    % \captionof{figure}{Problema \ref{p:P101}\label{f:P101}}
  \end{center}
  % \\[5pt]
  \begin{tabularx}{\textwidth}{XXXXX}
  $\square$ $x < \SI{-6}{\metre}$ & $\square$ $\SI{-2}{\metre} < x < 0$ & $\square$ $0 < x < \SI{1}{\metre}$ & $\square$ $\SI{1}{\metre} < x < \SI{2}{\metre}$ & $\square$ $ x > \SI{6}{\metre}$ \\
  \end{tabularx}

\

\textit{Resolución:}

Llamemos $q$ a la carga positiva que se agrega sobre el eje en la posición $x$. Si $0 < x < \SI{2}{\metre}$, la fuerza ejercida por $q_1$ será de repulsión hacia la derecha y la fuerza ejercida por $q_2$ será de atracción también hacia la derecha. Entonces la fuerza neta es hacia la derecha en cualquier posición entre las dos cargas.

\

Si $ x > \SI{2}{\metre}$ (a  la derecha de $q_2$), la fuerza que ejerce $q_2$ será de mayor módulo que la fuerza ejercida por $q_1$ porque $|q_2| > |q_1|$ y $q_2$ se encuentra más cerca de la tercera carga. Entonces no hay ninguna posición a la derecha de $q_2$ donde la fuerza neta sea hacia la derecha.

\ 

Si $x < 0$:
\begin{flalign*}
  \vec{F}_1 + \vec{F}_2 &= -\frac{kq_1q}{x^2} \hat{i} + \frac{kq_2q}{(\SI{2}{\metre}-x)^2}\hat{i}
\end{flalign*}
Por lo tanto, si buscamos cuando ese resultado es mayor a cero:
\begin{flalign*}
  kQq \left ( -\frac{1}{x^2} + \frac{2}{(\SI{2}{\metre}-x)^2} \right ) &> 0 
\end{flalign*}
La solución es $x < \SI{-4.83}{\metre}$. Entonces en $x < \SI{-6}{\metre}$ también se tiene una fuerza neta hacia la derecha.

\item
\begin{minipage}[t]{.6\textwidth}
    La figura muestra un hilo infinito con densidad de carga $\lambda = \SI{300}{\micro\coulomb/\metre}$, una carga puntual $q = \SI{20}{\micro\coulomb}$, y los puntos $A$ y $B$, todos sobre el mismo plano. Las distancias mostradas son $a = \SI{1}{\centi\metre}$ y $b = \SI{1.5}{\centi\metre}$. \textit{a}) Calcular el módulo del campo eléctrico total en el punto $A$. \textit{b}) Calcular la diferencia de potencial eléctrico entre los puntos $A$ y $B$.
   \end{minipage}
  %
  \begin{minipage}[t]{.4\textwidth}
  \strut\vspace*{-\baselineskip}
  \begin{center}
    \begin{tikzpicture}[scale=0.5]
      \draw (0,3.5) ellipse (0.2 and 0.1);
      \draw (0,-1.5) +(180:0.2 and 0.1) arc (180:360:0.2 and 0.1);
      \draw [black] (-0.2,-1.5)--(-0.2,3.5) node[pos=0.75, left] {$\lambda$};
      \draw [black] (0.2,-1.5)--(0.2,3.5);
      \fill [](-2,0) circle(5pt) node[left] {$q$};
      \draw [](2,2) circle(4pt) node[below] {$A$};
      \draw [](4,2) circle(4pt) node[below] {$B$};
      \draw [{latex}-{latex}] (-2,-0.4)--(0,-0.4) node[midway,below] {$a$};
      \draw [{latex}-{latex}] (0,2.4)--(2,2.4) node[midway,above] {$a$};
      \draw [{latex}-{latex}] (2,2.4)--(4,2.4) node[midway,above] {$a$};
      \draw [{latex}-{latex}] (4.6,2)--(4.6,0) node[midway,right] {$b$};
      \draw [dotted] (-1.8,0)--(4.6,0);
    \end{tikzpicture}
    % \captionof{figure}{\label{hilo}}
  \end{center}
  \end{minipage}

  \

\textit{Resolución:}
\textit{a}) El campo en $A$ es la suma del campo producido por el hilo y el del producido por la carga puntual.

\

Asumiendo eje $x$ positivo hacia la derecha y eje $y$ positivo hacia arriba, el campo del hilo en $A$ es:
\begin{flalign*}
  \vec{E}_h &= \frac{2k\lambda}{a}\hat{x}
\end{flalign*}
\begin{minipage}[t]{.5\textwidth}
El módulo del campo eléctrico de la carga puntual en $A$ es:
\begin{flalign*}
  |\vec{E}_q| &= \frac{kq}{d^2}
\end{flalign*}
El campo resultante es:
\begin{flalign*}
  \vec{E} &= (|\vec{E}_h|+|\vec{E}_q|\cos(\theta))\hat{x} + |\vec{E}_q|\sin(\theta)\hat{y} 
\end{flalign*}
Y el módulo simplemente se calcula como:
\begin{flalign*}
  |\vec{E}| &= \sqrt{E_x^2+E_y^2} 
\end{flalign*}

\end{minipage}
\begin{minipage}[t]{.5\textwidth}
  \strut\vspace*{-\baselineskip}
  \begin{center}
    \begin{tikzpicture}[scale=0.75]
      \draw (0,3.5) ellipse (0.2 and 0.1);
      \draw (0,-1.5) +(180:0.2 and 0.1) arc (180:360:0.2 and 0.1);
      \draw [black] (-0.2,-1.5)--(-0.2,3.5) node[pos=0.75, left] {$\lambda$};
      \draw [black] (0.2,-1.5)--(0.2,3.5);
      \fill [](-2,0) circle(5pt) node[left] {$q$};
      \draw [](2,2) circle(4pt) node[below] {$A$};
      \draw [](4,2) circle(4pt) node[below] {$B$};
      \draw [{latex}-{latex}] (-2,-0.4)--(0,-0.4) node[midway,below] {$a$};
      \draw [{latex}-{latex}] (0,2.4)--(2,2.4) node[midway,above] {$a$};
      \draw [{latex}-{latex}] (2,2.4)--(4,2.4) node[midway,above] {$a$};
      \draw [{latex}-{latex}] (4.6,2)--(4.6,0) node[midway,right] {$b$};
      \draw [dotted] (-1.8,0)--(4.6,0);
      \draw [blue, {latex}-{latex}, thick] (-2,0)--(2,2) node[pos=0.25,above] {$d$};
      \draw [blue, -{latex}, thick] (2,2)--(6,4) node[pos=0.75,above] {$\vec{E}_q$};
      \draw [blue, -{latex}, thick] (2,2)--(6,2) node[pos=0.75,above] {$\vec{E}_h$};
    \end{tikzpicture}
    % \captionof{figure}{\label{hilo}}
  \end{center}
\end{minipage}

\textit{b})
La diferencia de potencial entre $A$ y $B$ total es la suma de la diferencia de potencial producida por el campo del hilo más la producida por el campo de la carga puntual:
\begin{flalign*}
  V_B - V_A &= \left ( 2k\lambda\ln(r_a/r_b) \right ) + \left ( \frac{kq}{r_{qB}^2} - \frac{kq}{r_{qA}^2} \right )  
\end{flalign*}


\item Sabiendo que el potencial en un punto $P$ sobre el eje de simetría perpendicular a un disco plano, a una distancia $z$ del centro, si el radio del disco es $R = \SI{5.0}{\centi\metre}$ y la densidad superficial de carga es $\sigma = \SI{1.0}{\nano\coulomb/\metre}$, es:
\begin{flalign*}
  V_P &= \frac{\sigma}{2\varepsilon_o} \left ( \sqrt{R^2+z^2} - z \right )
\end{flalign*}
Elegir un valor de una carga puntual y encuentre una posición donde hay que ubicarla para que el potencial eléctrico total en $z = \SI{2.0}{\centi\metre}$ sea cero.

  \

\textit{Resolución:}\\
Si se agrega una carga puntual $q$, el potencial neto en el $z$ indicado es cero:
\begin{flalign*}
  V_\text{disco} + V_q &= 0\\
  \frac{\sigma}{2\varepsilon_o} \left ( \sqrt{R^2+z^2} - z \right ) + \frac{kq}{r_{qP}} &= 0\\
  \SI{1.91}{\volt} + \frac{kq}{r_{qP}} &= 0\\
  \frac{q}{r_{qP}} &= \frac{\SI{-1.91}{\volt}}{k} 
\end{flalign*}

Por supuesto que la carga $q$ debe ser negativa para que se cumpla lo pedido. Elegir un valor de $q$ y buscar la distancia al punto $P$ donde hay que ubicarla para que se cumpla $q/r_{qP} = \SI{-1.91}{\volt}/k$.

\item Cuando una resistencia de $\SI{12}{\ohm}$ se conecta a una batería de $\SI{12}{\volt}$, la potencia que se disipa en dicha resistencia es $\SI{10.22}{\watt}$. ¿Cuánto vale la resistencia interna de la batería?

\

\textit{Resolución:}\\

Conociendo la potencia disipada en la resistencia $R$, podemos encontrar la corriente en el circuito:
\begin{flalign*}
i &= \sqrt{\frac{P}{R}} 
\end{flalign*}
Luego:
\begin{flalign*}
  \SI{12}{\volt} -ir_\text{interna} - iR &= 0
\end{flalign*}
y se despeja el valor de la resistencia interna.

\item Para el circuito mostrado en la figura, encontrar el valor de $\varepsilon$ sabiendo que $V_B - V_A = \SI{2.00}{\volt}$.
\begin{center}
  \begin{circuitikz}[scale=0.8]
    \draw (0,0) to[R=$\SI{1.23}{\kilo\ohm}$] (4,0)
    (0,5.2) to[battery2, l=$\varepsilon$] (4,5.2)
    (0,3.2) -- (0.5,3.2) to[battery2, l=$\SI{12}{\volt}$, invert] (1.5,3.2) to[R=$\SI{0.700}{\kilo\ohm}$] (4,3.2)
    (0,0) -- (0,1.2) to[R=$\SI{2.25}{\kilo\ohm}$] (0,3.2) to[R=$\SI{11.5}{\kilo\ohm}$] (0,5.2)
    (4,5.2) to[R=$\SI{4.45}{\kilo\ohm}$] (4,3.2) -- (4,0) 
    (0,1.2) to[R=$\SI{1.23}{\kilo\ohm}$] (4,1.2)
    (0,0) -- (0,-0.5)
    (4,0) -- (4,-0.5);
    \fill (4,-0.5) circle (3pt);
    \fill (0,-0.5) circle (3pt);
    \draw (0,-0.5) node[right]{$A$};
    \draw (4,-0.5) node[right]{$B$};
  \end{circuitikz}
  % \captionof{figure}{Problema \ref{p:P6051}\label{f:P6051}}
\end{center}

\

\textit{Resolución:}\\
Las resistencias de $\SI{1.23}{\kilo\ohm}$ están conectadas en paralelo y la resistencia equivalente de ese paralelo es $R_p = \SI{615}{\ohm}$. A continuación se muestran las corrientes de mallas elegidas para los cálculos:

\begin{center}
  \begin{circuitikz}[scale=0.8]
    \draw (4,0) to[R=$R_p$] (0,0)
    (0,5.2) to[battery2, l=$\varepsilon$] (4,5.2)
    (0,3.2) -- (0.5,3.2) to[battery2, l=$\SI{12}{\volt}$, invert] (1.5,3.2) to[R=$R_3$] (4,3.2)
    (0,0) -- (0,1.2) to[R=$R_4$] (0,3.2) to[R=$R_1$] (0,5.2)
    (4,5.2) to[R=$R_2$] (4,3.2) -- (4,0) 
    (0,0) -- (0,-0.5)
    (4,0) -- (4,-0.5);
    \fill (4,-0.5) circle (3pt);
    \fill (0,-0.5) circle (3pt);
    \draw (0,-0.5) node[right]{$A$};
    \draw (4,-0.5) node[right]{$B$};
    \draw [blue, -{latex}, thick] (3,0.5)--(1,0.5) node[pos=0.5,above] {$i_2$};
    \draw [blue, -{latex}, thick] (0.4,0.6)--(0.4,2.6) node[pos=0.5,right] {$i_2$};
    \draw [blue, -{latex}, thick] (0.5,2.8)--(3,2.8) node[pos=0.5,below] {$i_2$};
    \draw [blue, -{latex}, thick] (3.6,2.6)--(3.6,0.6) node[pos=0.5,left] {$i_2$};
    \draw [blue, -{latex}, thick] (1.6,4.8)--(3.2,4.8) node[pos=0.5,below] {$i_1$};
    \draw [blue, -{latex}, thick] (3.6,4.8)--(3.6,3.4) node[pos=0.5,left] {$i_1$};
    \draw [blue, -{latex}, thick] (3,3.4)--(1,3.4) node[pos=0.5,above] {$i_1$};
    \draw [blue, -{latex}, thick] (0.6,3.4)--(0.6,4.8) node[pos=0.95,right] {$i_1$};

  \end{circuitikz}
  % \captionof{figure}{Problema \ref{p:P6051}\label{f:P6051}}
\end{center}

La ecuación de la malla 1 es:
\begin{flalign*}
    -\varepsilon - \SI{12}{\volt} -(R_1+R_2+R_3)i_1 + R_3i_2 &= 0 
\end{flalign*}
Y para la malla 2 es:
\begin{flalign*}
  \SI{12}{\volt} +R_3i_1 -(R_3+R_4+R_p)i_2 &= 0 
\end{flalign*}

Sabiendo la diferencia de potencial entre $A$ y $B$ se puede encontrar la corriente de malla $i_2$, siendo la corriente que circula a través de $R_p$:
\begin{flalign*}
  i_2 &= \frac{\SI{2}{\volt}}{R_p}
\end{flalign*}

Reemplazando este valor en la ecuación de la malla 2 se puede encontrar el valor de $i_1$, y luego reemplazando ambas corrientes en la ecuación de la malla 1 se puede encontrar la respuesta.


\end{enumerate}

\end{document}
